% ============================================================
%  Section 5 – Interprétation et discussion
% ============================================================
\section{Interprétation et discussion}

\subsection{Interprétation des résultats}

\textbf{T1.} Dans notre population, 80,9\,\% des répondant·es dorment moins de
7~heures avant un DE. Le test montre que cet écart à 50\,\% est statistiquement
significatif ($z = 5{,}09$) : le manque de sommeil avant les examens est une
tendance forte et non attribuable au hasard de l'échantillonnage.

\textbf{T2.} La proportion d'étudiant·es assistant à tous les CMs est de
$\hat{p} = 42{,}6\,\%$. On ne peut cependant pas affirmer, au risque~5\,\%,
qu'elle dépasse $1/3$ ($z = 1{,}63 < 1{,}6449$). Le résultat est très proche
du seuil, ce qui le rend fragile à une légère variation de l'échantillon.

\textbf{T3.} La moyenne observée d'annales travaillées ($\bar{x} \approx 2{,}19$)
n'est pas significativement différente de~2 ($t = 1{,}78 < 1{,}9960$). L'IC~95\,\%
$[1{,}977\,;\,2{,}406]$ contient la valeur~2, ce qui confirme ce résultat.

\textbf{T4.} La distribution de V12 n'est pas uniforme ($\chi^2 = 43{,}03 \gg 7{,}815$) :
58,2\,\% des répondant·es se situent dans la tranche 12--14, ce qui traduit une
certaine homogénéité du niveau académique initial dans notre groupe.

\textbf{T5.} Les conditions du $\chi^2$ d'indépendance entre V2 et V14 ne sont
pas satisfaites. Aucune conclusion statistique n'est possible sur ce lien.

\textbf{T6.} On ne met pas en évidence de dépendance significative entre la
fréquence de participation extra-scolaire et la moyenne au S1
($\chi^2 = 4{,}375 < 9{,}488$) : les données disponibles ne permettent pas de
conclure à une association entre ces deux variables.

\subsection{Association statistique et causalité}

Il est essentiel de distinguer association statistique et relation causale.
Dans notre étude, même lorsqu'un test rejette $H_0$, cela ne prouve pas qu'une
variable en cause une autre. Par exemple, la proportion élevée d'étudiant·es
dormant peu (T1) ne signifie pas que le manque de sommeil cause de mauvais
résultats : d'autres facteurs (charge de travail, stress, organisation) peuvent
expliquer simultanément les deux phénomènes.

De même, l'absence d'association détectée ne garantit pas l'absence de lien
dans la réalité. Notre protocole est observationnel ; les réponses sont
déclaratives et non mesurées objectivement. Toute interprétation doit donc
rester prudente.

\subsection{Limites de l'étude}

\begin{itemize}
  \item \textbf{Effectif limité} ($n = 68$) : certains tests du $\chi^2$ ne
        peuvent pas être appliqués (T5, conditions non satisfaites).
  \item \textbf{Représentativité} : participation volontaire sur 4~classes ;
        nos résultats ne sont pas généralisables à l'ensemble de la P1 EFREI.
  \item \textbf{Auto-déclaration} : les réponses sont subjectives et exposées
        au biais de désirabilité sociale (assiduité ou travail personnel
        potentiellement surestimés).
  \item \textbf{Variables ordinales} parfois traitées comme quantitatives
        (V2) : approximation nécessaire pour les tests de moyenne.
  \item \textbf{Formulation des questions} : certaines modalités du sondage
        restent imprécises (ex.\ : \og{}4+\fg{} pour V2 compresse les valeurs
        élevées), ce qui peut biaiser les comparaisons.
  \item \textbf{Périmètre de T6} : la question des permanences pédagogiques
        n'a pas pu être testée directement en raison de contraintes d'effectifs ;
        T6 porte sur les activités extra-scolaires à titre de substitut.
\end{itemize}

\section*{Conclusion}

Notre étude montre qu'une grande majorité des étudiant·es interrogé·es dorment
moins de 7~heures avant un DE (T1 rejeté) et que les niveaux de bac sont
concentrés sur la tranche 12--14 (T4 rejeté). En revanche, on ne peut pas
conclure à une différence significative du nombre moyen d'annales travaillées,
ni à une dépendance entre la participation extra-scolaire et la moyenne au S1.

Ces résultats permettent d'identifier certaines tendances dans nos données, mais
notre protocole observationnel ne permet pas d'établir de relations causales.
Pour aller plus loin, un effectif plus large et des données objectives
(notes réelles, présences enregistrées) renforceraient la portée des conclusions.
